\begin{resumo}[Abstract]
 \begin{otherlanguage*}{english}
   Government data platforms already consolidate the execution history of public policies, but they rarely have the infrastructure to put machine learning models into production reliably. This work investigates how to structure an MLOps architecture integrated into a government data platform organized in medallion layers, so that the results of predictive models are reproducible, traceable and comparable to the current decision rule without temporal information leakage. The proposed architecture extends the existing Data Lakehouse instead of building a separate platform: the panel and the label live in the silver layer; features computed as of each observation date, scores and monitoring metrics live in the gold layer; and scores re-enter through the bronze layer like any other source, under the same tests and the same catalog. The architecture is defined by tool-independent roles: the transformation layer performs all computation on data, the modeling environment only trains, evaluates and scores, the orchestrator coordinates the steps, and each training run freezes a versioned dataset recorded with the experiment. The same roles also accommodate applications built on language models, which is discussed but not evaluated. The research follows Design Science Research, and the architecture is instantiated in the prediction of construction stoppages in the Minha Casa, Minha Vida housing program on the Gov Hub Cidades platform of the Brazilian Ministry of Cities, formulated as discrete-time survival analysis and compared with the flag-based rule currently used by inspectors. In this first stage, the requirements and the architecture design were defined, and the data survey showed that source changes along the historical series limit the period in which the stoppage label is reliable; the recovery of a legacy series found labels for 2015 to 2018 as well, which enter the training base under explicit acceptance criteria. In the second stage, the architecture will be implemented, demonstrated by replaying the monthly operation over the historical series, and evaluated by eight criteria with success conditions defined in advance: adherence to MLOps principles, production readiness, reproducibility, absence of temporal leakage, comparison with the current rule, fairness across subgroups, traceability and extensibility.

   \vspace{\onelineskip}

   \noindent
   \textbf{Keywords}: MLOps. government data. data lakehouse. medallion architecture. survival analysis.
 \end{otherlanguage*}
\end{resumo}
